\section{Related Work}
\label{sec:related}

\noindent\textbf{Device fingerprinting in cyber-physical systems.} Formby et~al.~\cite{formby2016}
fingerprint ICS devices from cross-layer response timing and physical-operation timing, distinguishing
device types and instances without decoding payloads, and GTID~\cite{gtid2015} identifies devices and
device types from wire-side timing signatures. Physical-layer signatures fingerprint devices even more
directly, from clock skew observed remotely~\cite{kohno2005} to the process-actuator physics of a
cyber-physical device~\cite{gu2018}. This line establishes the threat the defense addresses:
observable timing and shape features identify an outstation. It is an attacker capability rather than a
defense. The defense differs in objective; we suppress the master-facing CLRT and segment-shape features
these techniques exploit, for one evaluated device, rather than proposing a new fingerprint.

\noindent\textbf{Network-timing traffic analysis and its defenses.} Inter-event timing alone leaks
secret content, from keystroke intervals on SSH~\cite{song2001} to modern website-fingerprinting
attacks that recover a browsed page from packet sizes and timing with a deep
classifier~\cite{sirinam2018}, although such accuracy is contested at realistic open-world
scale~\cite{panchenko2016}. The corresponding defenses shape the traffic: adaptive padding inserts
dummy packets to break the inter-packet timing pattern~\cite{juarez2016}, and comprehensive
side-channel mitigation shapes a tenant's traffic to a schedule that leaks no secret~\cite{pacer2022}.
These defenses target browser traffic over an anonymity network, or cloud tenants, and assume a
cooperating end host or hypervisor that can add cover traffic. Unlike these host-cooperative defenses, our mechanism normalizes the ACK-to-response timing of a
protocol-bound DNP3 exchange in the network, with no endpoint change and no added application bytes,
and it fixes a specific cross-layer interval rather than obscuring an inter-packet distribution.

\noindent\textbf{Packet-size padding, morphing, and segmentation.} Size-based countermeasures reshape
the length signature of a flow. Traffic morphing optimally transforms one class's packet-size
distribution into another's~\cite{wright2009}, and later analysis shows that per-packet padding and
fragmentation leave count and size features that a classifier still exploits~\cite{dyer2012}.
Constant-rate padding such as Tamaraw bounds leakage by fixing the packet size and interval, at a cost
acceptable for one flow but prohibitive across a device fleet~\cite{tamaraw2014}, while burst molding
collides distinct activities onto one pattern with lower overhead~\cite{walkietalkie2017}. The same
size-and-rate channel identifies devices even under encryption, where traffic shaping is the standard
countermeasure~\cite{apthorpe2019}, and shaping can be made differentially private to bound the
residual leak~\cite{netshaper2024}. These lines change or pad bytes and target the end host. The defense differs in mechanism; it neither pads nor
rewrites DNP3 bytes, but re-segments the response only at existing CRC-block boundaries, so the
reassembled message and every per-block CRC are preserved while the observed segment vector is fixed.

\noindent\textbf{Anti-reconnaissance and moving-target defense for power grids.} A body of work
disrupts grid reconnaissance rather than detecting its use. DefRec virtualizes physical functions to
present deceptive content and connectivity to an attacker~\cite{defrec2020}, RAINCOAT randomizes data
acquisition and injects decoy measurements to mislead attack strategy~\cite{raincoat2019}, and a
companion testbed evaluates such anti-reconnaissance approaches on grid cyber-physical
infrastructure~\cite{lintestbed2020}. A parallel detection line adapts specification-based intrusion
detection to DNP3~\cite{linbro2013}. In-network obfuscation has denied reconnaissance before: NetHide
rewrites data-plane topology to blunt link-flooding attacks while keeping path tracing
usable~\cite{nethide2018}. These defenses operate at the content, connectivity, topology, or detection
layer. Our defense instead operates at a different observation point: it obfuscates the wire-visible size and timing of the
transport exchange itself, in the network and below the application semantics those defenses reshape,
and it complements rather than replaces them.

\noindent\textbf{Timing and shaping on programmable switches.} Programmable data planes make in-network
traffic shaping practical. The P4 model exposes a programmable parser and match-action
pipeline~\cite{p4ccr2014}, and push-in-first-out queues enable programmable packet scheduling at line
rate~\cite{pifo2016}. PIFO is a hardware abstraction; SP-PIFO shows that rank-ordered scheduling can be
approximated with a few strict-priority queues on today's switches~\cite{sppifo2020}, the primitive
class our deadline-driven response scheduler is built on. Ditto pads packets and injects chaff on a
programmable switch to make WAN traffic
follow a deterministic pattern independent of the payload, at line rate and without end-host
changes~\cite{ditto2022}. Ditto is the closest in-network precedent. It adds bytes and packets to reach
a fixed volume pattern, which is appropriate for opaque WAN traffic but not for a protocol whose bytes
and CRCs the receiver validates. Recent systems bring traffic-analysis defense fully into the
programmable data plane: Securitas obfuscates flows against traffic analysis with flexible in-network
shaping~\cite{securitas2026}, Minos mounts a lightweight, dynamic defense against traffic analysis on
programmable switches~\cite{minos2025}, and, closest to our setting, Hu and Lin enhance industrial
network-protocol security with programmable switches~\cite{hulin2023}. The defense differs in two
ways from all of these: it is byte-preserving and carves only at existing DNP3 boundaries, so the
outstation's protocol remains valid, and it adds a control-command mode that hides operation timing
without exposing the hold to a passive observer.
